\scanpage{139}

(30)\qquad $\forall E\in\mathcal E_0$, \deleted{\unsure{une}}la donnée de la structure de $n$-polygone combinatoire, ayant $E$ comme ens. de ses sommets, i.e. \inserted{de} $A_E\subset\mathfrak P_2(E)$

puisqu'on récupère $A_0$ à partir de (29) (30) par
\[
A_0=\{a\in\mathfrak P_2(S)\mid\exists E\in\mathcal E_0\ \text{t.q. }a\in A_E\}.\tag*{(31)}
\]
À vrai dire, \deleted{pour} dans le cas où l'axiome (To 1) est satisfait, il existe exactement \uline{deux} $s\in S$ tels que $\mathrm{Et}_{\deleted{\unsure{\lambda}}A_0}(s)$ contienne $a$ (donc admette $a$ comme arête) et sauf dans le cas \unsure{octaèdral|octaédral}, ces étoiles associées sont distinctes, donc sauf dans le cas de l'octaèdre, il y aura dans (31) exactement deux $E\in\mathcal E_0$ telles que $E\supset a$, i.e $a\in A_E$. (Dans le cas de l'octaèdre, par contre, le $E$ en question est unique.).

Pour déterminer \deleted{\unsure{une}} \unsure{structure|structures} $A_\lambda$ $(\lambda\in Z)$, comme \inserted{l'ens $\mathcal E_\lambda$ de} ses étoiles est égal à $\mathcal E_0$, il suffit de savoir décrire, \deleted{$\unsure{E}$} pour $\forall E\in\mathcal E_0$, la \deleted{\unsure{struc}}\inserted{structure} $n$-polygonale $A_{E,\lambda}\in\mathfrak P_2(E)$ induite sur $E$ par $\Gamma_\lambda=(S,A_\lambda)$. Notons maintenant le fait mirifique

(32)\qquad Soit $E$ un ens. \inserted{fini,} de cardinal $n\geqslant3$. Alors le groupe $Z'_n=(\mathbb Z/n\mathbb Z)^*/\{\pm1\}$ opère canoniquement sur l'ens. \inserted{$\mathrm{Pol}(E)$} des structures polygônales sur $E$ comme ens de sommets :

% slanted
\margin{On voit de suite que cette opération est fidèle.}

% archive sheet number: 138
